\chapter{Convex Functions on Riemannian Manifolds}
\label{app:chow-iii-convex-functions-riemannian-manifolds}

\providecommand{\Lip}{\operatorname{Lip}}
\providecommand{\dil}{\operatorname{dil}}
\providecommand{\R}{\mathbb R}
\providecommand{\M}{\mathcal M}
\providecommand{\C}{\mathcal C}
\providecommand{\T}{\widetilde T}
\providecommand{\Xzero}{\mathfrak X_0}
\providecommand{\Vzero}{\mathcal V_0}
\providecommand{\dC}{d_C}

This appendix develops the elementary convex analysis of locally convex subsets
of Riemannian manifolds and the distance-nonincreasing retraction used later.

\section{Euclidean convexity}

\begin{proposition}[Rademacher theorem on manifolds]
\label{prop:chow-iii-h2-rademacher}
\source{chow-et-al-ricci-flow-part-iii:page-0435}
\dcref{apph:H.2}\group{Euclidean convexity}
If $f:\M\to\R$ is locally Lipschitz, then $f$ is differentiable almost everywhere.
In particular, distance functions are differentiable off a measure-zero cut locus.
\end{proposition}

\begin{lemma}[Fundamental theorem for Lipschitz functions]
\label{lem:chow-iii-h3-lipschitz-ftc}
\source{chow-et-al-ricci-flow-part-iii:page-0435}
\dcref{apph:H.3}\group{Euclidean convexity}
If $f:[a,b]\to\R$ is Lipschitz, then
\[
f(x)=f(a)+\int_a^x f'(s)\,ds
\]
for every $x\in[a,b]$, with the derivative defined almost everywhere.
\end{lemma}

\begin{remark}[Extended-valued convex extension]
\label{rem:chow-iii-h4-extended-convex-extension}
\source{chow-et-al-ricci-flow-part-iii:page-0436}
\dcref{apph:H.4}\group{Euclidean convexity}
If $C\subset\R^n$ is convex and $f:C\to\R$ is convex, its extension by $+\infty$
outside $C$ is convex on all of $\R^n$.
\end{remark}

\begin{proposition}[Interior regularity of Euclidean convex functions]
\label{prop:chow-iii-h5-convex-lipschitz}
\source{chow-et-al-ricci-flow-part-iii:page-0436}
\dcref{apph:H.5}\group{Euclidean convexity}
Every finite-valued convex function on a locally convex $C\subset\R^n$ is locally
Lipschitz, and hence differentiable almost everywhere, on $\operatorname{int}(C)$.
\end{proposition}

\begin{lemma}[Pointwise limits of convex functions]
\label{lem:chow-iii-h6-convex-limits}
\source{chow-et-al-ricci-flow-part-iii:page-0436}
\dcref{apph:H.6}\group{Euclidean convexity}
If finite-or-infinite-valued convex functions $f_k:\R^n\to\R\cup\{\infty\}$ converge
pointwise to $f_\infty$, then $f_\infty$ is convex and the convergence is uniform on
compact subsets where the limit is finite.
\end{lemma}

\begin{definition}[Twice differentiability in the sense of Stolz]
\label{def:chow-iii-h7-stolz-differentiability}
\source{chow-et-al-ricci-flow-part-iii:page-0437}
\dcref{apph:H.7}\group{Euclidean convexity}
For an open $U\subset\R^n$, a continuous $f$ is twice differentiable in the sense of
Stolz at $x$ if there are a vector $W$ and symmetric matrix $M$ such that
\[
f(x+y)=f(x)+W\cdot y+M(y,y)+o(|y|^2).
\]
\end{definition}

\begin{lemma}[Nearest-point projection inequality]
\label{lem:chow-iii-h8-projection-inequality}
\source{chow-et-al-ricci-flow-part-iii:page-0437}
\dcref{apph:H.8}\group{Euclidean convexity}
For a closed convex $K\subset\R^n$, its nearest-point projection satisfies
\[
\langle x-\pi_K(x),z-\pi_K(x)\rangle\le0
\qquad(x\in\R^n,\ z\in K).
\]
\end{lemma}

\begin{proposition}[Projection is distance nonincreasing]
\label{prop:chow-iii-h9-projection-contraction}
\source{chow-et-al-ricci-flow-part-iii:page-0437}
\uses{lem:chow-iii-h8-projection-inequality}
\dcref{apph:H.9}\group{Euclidean convexity}

For a closed convex $K\subset\R^n$,
\[
|\pi_K(x)-\pi_K(y)|\le |x-y|.
\]
\end{proposition}
\begin{proof}
\sketch
Apply the projection inequality twice and add the resulting inequalities.
\end{proof}

\begin{corollary}[Projected paths do not increase length]
\label{cor:chow-iii-h10-projection-length}
\source{chow-et-al-ricci-flow-part-iii:page-0438}
\uses{prop:chow-iii-h9-projection-contraction}
\dcref{apph:H.10}\group{Euclidean convexity}

For every rectifiable path $\gamma$ in $\R^n$,
\[
\operatorname{L}(\pi_K\circ\gamma)\le \operatorname{L}(\gamma).
\]
\end{corollary}

\section{Locally convex subsets}

\begin{definition}[Interior tangent cone]
\label{def:chow-iii-h11-interior-tangent-cone}
\source{chow-et-al-ricci-flow-part-iii:page-0439}
\dcref{apph:H.11}\group{Locally convex subsets}
For $S\subset\M$ and $p\in S$, define $\T_pS$ as the set of vectors $V$ for which
$\exp_p(sV)\in\operatorname{int}(S)$ for all sufficiently small $s>0$.
\end{definition}

\begin{remark}[Tangent-cone realization]
\label{rem:chow-iii-h12-tangent-cone-realization}
\source{chow-et-al-ricci-flow-part-iii:page-0439}
\dcref{apph:H.12}\group{Locally convex subsets}
If a metric tangent cone exists and embeds in $T_p\M$, it is expected to agree with
the closure of the interior tangent cone.
\end{remark}

\begin{lemma}[Openness of the interior tangent cone]
\label{lem:chow-iii-h13-tangent-cone-open}
\source{chow-et-al-ricci-flow-part-iii:page-0439}
\dcref{apph:H.13}\group{Locally convex subsets}
If $S$ is convex, then $\T_pS$ is open in $T_p\M$.
\end{lemma}

\begin{remark}[Variants of geodesic convexity]
\label{rem:chow-iii-h14-convexity-variants}
\source{chow-et-al-ricci-flow-part-iii:page-0440}
\dcref{apph:H.14}\group{Locally convex subsets}
A subset is convex when every minimizing ambient geodesic between two of its points
stays in the subset; it is locally convex when this holds in a neighborhood of each
point.  Intersections of convex sets are convex, while locally convex sets need not
be connected.
\end{remark}

\begin{lemma}[Minimal geodesics in a closed locally convex set]
\label{lem:chow-iii-h15-minimal-geodesics-locally-convex}
\source{chow-et-al-ricci-flow-part-iii:page-0441,chow-et-al-ricci-flow-part-iii:page-0442,chow-et-al-ricci-flow-part-iii:page-0443}
\dcref{apph:H.15}\group{Locally convex subsets}
If $C$ is connected, closed, and locally convex in a complete Riemannian manifold,
then every $x,y\in C$ are joined by a smooth ambient geodesic in $C$ of length
\[
d_C(x,y)=\inf\{\operatorname{L}(\beta):\beta\subset C,\ \beta(0)=x,\ \beta(1)=y\}.
\]
\end{lemma}

\begin{remark}[Locally convex sets as length spaces]
\label{rem:chow-iii-h16-locally-convex-length-space}
\source{chow-et-al-ricci-flow-part-iii:page-0443}
\dcref{apph:H.16}\group{Locally convex subsets}
The preceding minimizing-geodesic statement can alternatively be obtained by
minimizing among piecewise-short broken geodesics.
\end{remark}

\begin{proposition}[Aleksandrov structure of locally convex subsets]
\label{prop:chow-iii-h17-aleksandrov-subset}
\source{chow-et-al-ricci-flow-part-iii:page-0443}
\uses{lem:chow-iii-h15-minimal-geodesics-locally-convex}
\dcref{apph:H.17}\group{Locally convex subsets}

If $\sec_{\M}\ge k$ and $C$ is connected, closed, and locally convex, then the
length space $C$ is an Aleksandrov space with curvature bounded below by $k$.
\end{proposition}

\begin{lemma}[Infinitesimal convexity]
\label{lem:chow-iii-h18-infinitesimal-convexity}
\source{chow-et-al-ricci-flow-part-iii:page-0444,chow-et-al-ricci-flow-part-iii:page-0445}
\dcref{apph:H.18}\group{Locally convex subsets}
If $V_1,V_2$ are unit vectors whose short radial geodesics lie in $\operatorname{int}(C)$,
then all sufficiently short radial geodesics in directions between them lie in
$\operatorname{int}(C)$.  The same holds when the first radial geodesic lies in $C$
and the second in its interior, provided the interpolation parameter is positive.
\end{lemma}

\begin{corollary}[Convexity of the interior tangent cone]
\label{cor:chow-iii-h19-interior-tangent-cone-convex}
\source{chow-et-al-ricci-flow-part-iii:page-0445}
\uses{lem:chow-iii-h18-infinitesimal-convexity}
\dcref{apph:H.19}\group{Locally convex subsets}

For a locally convex $C$, the cone $\T_pC$ is convex.
\end{corollary}

\section{Convex functions and directional derivatives}

\begin{definition}[Convex function on a locally convex set]
\label{def:chow-iii-h20-convex-function}
\source{chow-et-al-ricci-flow-part-iii:page-0446}
\dcref{apph:H.20}\group{Convex functions}
A function $f:C\to\R$ is convex when $f\circ\gamma$ is convex on every constant-speed
geodesic $\gamma$ contained in $C$; it is concave when $-f$ is convex.
\end{definition}

\begin{lemma}[Convex functions are locally Lipschitz]
\label{lem:chow-iii-h21-convex-locally-lipschitz}
\source{chow-et-al-ricci-flow-part-iii:page-0446,chow-et-al-ricci-flow-part-iii:page-0447,chow-et-al-ricci-flow-part-iii:page-0448}
\uses{def:chow-iii-h20-convex-function}
\dcref{apph:H.21}\group{Convex functions}

Every convex function on a connected locally convex set is locally Lipschitz, and
therefore continuous, on the interior.
\end{lemma}

\begin{definition}[Directional derivative]
\label{def:chow-iii-h22-directional-derivative}
\source{chow-et-al-ricci-flow-part-iii:page-0449}
\dcref{apph:H.22}\group{Directional derivatives}
For $V\in\T_pS$, define
\[
(D_Vf)(p)=\lim_{s\to0^+}\frac{f(\exp_p(sV))-f(p)}{s}
\]
when the limit exists.  The difference quotient is denoted $J_s(V)$.
\end{definition}

\begin{lemma}[Homogeneity of directional derivatives]
\label{lem:chow-iii-h23-directional-homogeneity}
\source{chow-et-al-ricci-flow-part-iii:page-0449}
\uses{def:chow-iii-h22-directional-derivative}
\dcref{apph:H.23}\group{Directional derivatives}

For $r>0$, $D_{rV}f=rD_Vf$ whenever the directional derivatives are defined.
\end{lemma}

\begin{remark}[Euclidean difference quotients]
\label{rem:chow-iii-h24-euclidean-difference-quotient}
\source{chow-et-al-ricci-flow-part-iii:page-0449}
\dcref{apph:H.24}\group{Directional derivatives}
In Euclidean space $J_s(V)=(f(p+sV)-f(p))/s$.
\end{remark}

\begin{lemma}[Directional derivatives of convex functions]
\label{lem:chow-iii-h25-directional-properties}
\source{chow-et-al-ricci-flow-part-iii:page-0449,chow-et-al-ricci-flow-part-iii:page-0450,chow-et-al-ricci-flow-part-iii:page-0451}
\uses{def:chow-iii-h22-directional-derivative}
\dcref{apph:H.25}\group{Directional derivatives}

For convex $f$, $J_{s_1}(V)\le J_{s_2}(V)$ when $0<s_1\le s_2$, so $D_Vf$ exists
in $[-\infty,\infty)$.  If $f$ is locally $L$-Lipschitz, then
\[
|J_s(V)|\le L|V|,\qquad |D_Vf|\le L|V|,
\]
and if both $V$ and $-V$ are admissible then $-D_{-V}f\le D_Vf$.
\end{lemma}

\begin{lemma}[Convexity and semicontinuity of $Df$]
\label{lem:chow-iii-h26-directional-convexity-usc}
\source{chow-et-al-ricci-flow-part-iii:page-0451,chow-et-al-ricci-flow-part-iii:page-0452,chow-et-al-ricci-flow-part-iii:page-0453}
\uses{lem:chow-iii-h25-directional-properties}
\dcref{apph:H.26}\group{Directional derivatives}

For convex $f$, $V\mapsto D_Vf(p)$ is convex on $\T_pC$ (in Euclidean space directly,
and on a locally Lipschitz Riemannian neighborhood by blow-up).  The function
$(p,V)\mapsto D_Vf(p)$ is upper semicontinuous.
\end{lemma}

\begin{lemma}[Sublinearity of the directional derivative]
\label{lem:chow-iii-h27-directional-sublinear}
\source{chow-et-al-ricci-flow-part-iii:page-0453}
\uses{lem:chow-iii-h23-directional-homogeneity, lem:chow-iii-h26-directional-convexity-usc}
\dcref{apph:H.27}\group{Directional derivatives}

If $f$ is locally $L$-Lipschitz near $p$, then
\[
|D_Wf(p)-D_Vf(p)|\le L|W-V|.
\]
In the interior, $V\mapsto D_Vf(p)$ is a finite sublinear function.
\end{lemma}

\begin{remark}[Concave directional derivatives]
\label{rem:chow-iii-h28-concave-directional-derivatives}
\source{chow-et-al-ricci-flow-part-iii:page-0453}
\dcref{apph:H.28}\group{Directional derivatives}
For a concave function in the interior, directional derivatives are finite, positively
homogeneous, concave, and satisfy $-D_{-V}f\ge D_Vf$.
\end{remark}

\section{Generalized gradients}

\begin{definition}[Generalized gradient]
\label{def:chow-iii-h29-generalized-gradient}
\source{chow-et-al-ricci-flow-part-iii:page-0454}
\uses{lem:chow-iii-h27-directional-sublinear}
\dcref{apph:H.29}\group{Generalized gradients}

For a locally Lipschitz convex $f$, choose a unit minimizer $W$ of $D_Vf(p)$ on
$\overline{\T_pC}\cap S_p^{n-1}$ and set
\[
\nabla f(p)=|D_Wf(p)|W.
\]
Thus $|\nabla f(p)|=|D_{\min}f(p)|$.
\end{definition}

\begin{lemma}[Elementary generalized-gradient properties]
\label{lem:chow-iii-h30-generalized-gradient-properties}
\source{chow-et-al-ricci-flow-part-iii:page-0454,chow-et-al-ricci-flow-part-iii:page-0455}
\uses{def:chow-iii-h29-generalized-gradient}
\dcref{apph:H.30}\group{Generalized gradients}

At differentiability points the generalized gradient is the negative of the ordinary
gradient.  It need not be unique; moreover
\[
D_{\nabla f(p)}f(p)=\operatorname{sign}(D_{\min}f(p))|\nabla f(p)|^2,
\qquad |\nabla f(p)|\le\Lip(f).
\]
\end{lemma}

\begin{lemma}[Interior critical points are minima]
\label{lem:chow-iii-h31-critical-local-minimum}
\source{chow-et-al-ricci-flow-part-iii:page-0455}
\uses{def:chow-iii-h29-generalized-gradient}
\dcref{apph:H.31}\group{Generalized gradients}

If $p\in\operatorname{int}(C)$ and $D_{\min}f(p)\ge0$, then $p$ is a local minimum
of the convex function $f$.
\end{lemma}

\begin{definition}[Convex functions vanishing on the boundary]
\label{def:chow-iii-h32-Xzero}
\source{chow-et-al-ricci-flow-part-iii:page-0456}
\dcref{apph:H.32}\group{Generalized gradients}
The class $\Xzero(C)$ consists of locally Lipschitz convex functions $f:C\to\R$ on
closed connected locally convex sets with nonempty boundary and interior, satisfying
$f|_{\partial C}=0$.
\end{definition}

\begin{lemma}[Boundary extension and convexity of $Df$]
\label{lem:chow-iii-h33-boundary-directional-extension}
\source{chow-et-al-ricci-flow-part-iii:page-0456,chow-et-al-ricci-flow-part-iii:page-0457}
\uses{def:chow-iii-h32-Xzero, lem:chow-iii-h26-directional-convexity-usc}
\dcref{apph:H.33}\group{Generalized gradients}

For $f\in\Xzero(C)$ and $p\in\partial C$, $D_Vf(p)$ tends to zero as $V$ approaches
the boundary of $\T_pC\cap S_p^{n-1}$; its continuous extension to the closed cone is
convex and locally Lipschitz.
\end{lemma}

\begin{example}[The square distance-to-boundary function]
\label{ex:chow-iii-h34-square-convex-function}
\source{chow-et-al-ricci-flow-part-iii:page-0457}
\dcref{apph:H.34}\group{Generalized gradients}
On $C=[-1,1]^2$, the function $f(x,y)=-d((x,y),\partial C)=\max\{|x|,|y|\}-1$
belongs to $\Xzero(C)$ and is nonsmooth on $\{|x|=|y|\}$.
\end{example}

\begin{remark}[Generalized gradient in the square example]
\label{rem:chow-iii-h35-square-gradient}
\source{chow-et-al-ricci-flow-part-iii:page-0457,chow-et-al-ricci-flow-part-iii:page-0458}
\dcref{apph:H.35}\group{Generalized gradients}
In the square example the generalized gradient takes the four axial values off the
diagonals and the four half-diagonal values on them; neither $\nabla f$ nor its norm
is continuous across the diagonals.
\end{remark}

\begin{lemma}[Uniqueness of the generalized gradient]
\label{lem:chow-iii-h36-unique-generalized-gradient}
\source{chow-et-al-ricci-flow-part-iii:page-0458,chow-et-al-ricci-flow-part-iii:page-0459,chow-et-al-ricci-flow-part-iii:page-0460}
\uses{lem:chow-iii-h15-minimal-geodesics-locally-convex, lem:chow-iii-h33-boundary-directional-extension}
\dcref{apph:H.36}\group{Generalized gradients}

For $f\in\Xzero(C)$, every $p\in C$ has a nonempty interior tangent cone and a
generalized gradient in its closure.  If $f_{\inf}<0$ and $f(p)>f_{\inf}$, the negative
minimum of $D_Vf(p)$ on $\T_pC\cap S_p^{n-1}$ is attained at a unique vector, so
$\nabla f(p)$ is unique and nonzero.
\end{lemma}

\begin{lemma}[Directional derivatives and the generalized gradient]
\label{lem:chow-iii-h37-directional-gradient-inequality}
\source{chow-et-al-ricci-flow-part-iii:page-0460,chow-et-al-ricci-flow-part-iii:page-0461}
\uses{lem:chow-iii-h36-unique-generalized-gradient}
\dcref{apph:H.37}\group{Generalized gradients}

For $f\in\Xzero(C)$ with $f_{\inf}<0$,
\[
D_Vf(p)\ge-\langle\nabla f(p),V\rangle
\]
for every $V$ in the tangent space at a minimum point and for every $V\in\T_pC$
at a non-minimum point.
\end{lemma}

\begin{lemma}[Acute-angle property]
\label{lem:chow-iii-h38-acute-angle}
\source{chow-et-al-ricci-flow-part-iii:page-0461}
\uses{lem:chow-iii-h37-directional-gradient-inequality}
\dcref{apph:H.38}\group{Generalized gradients}

If $f(x)>f_{\inf}$, $f(x)\ge f(y)$, and a minimizing geodesic from $x$ to $y$ has
initial velocity $\dot\gamma(0)$, then
\[
\angle(\nabla f(x),\dot\gamma(0))\le\frac\pi2.
\]
\end{lemma}

\begin{lemma}[Nonnegative Hessian in the support sense]
\label{lem:chow-iii-h39-convex-support-hessian}
\source{chow-et-al-ricci-flow-part-iii:page-0462}
\uses{lem:chow-iii-h37-directional-gradient-inequality}
\dcref{apph:H.39}\group{Generalized gradients}

At every interior point of a connected locally convex set, a convex function has
$\operatorname{Hess}_{\mathrm{sup}}f\ge0$ and is twice differentiable in the sense of
Stolz almost everywhere.
\end{lemma}

\section{Concave integral curves and retraction}

\begin{definition}[Right tangent vector]
\label{def:chow-iii-h40-right-tangent}
\source{chow-et-al-ricci-flow-part-iii:page-0463}
\dcref{apph:H.40}\group{Integral curves}
A vector $\dot\gamma_+(s)\in T_{\gamma(s)}\M$ is the right tangent of a continuous path
when every smooth $h$ satisfies $(h\circ\gamma)'_+(s)=\dot\gamma_+(s)(h)$.
\end{definition}

\begin{lemma}[Geodesic characterization of the right tangent]
\label{lem:chow-iii-h41-right-tangent-geodesics}
\source{chow-et-al-ricci-flow-part-iii:page-0464}
\uses{def:chow-iii-h40-right-tangent}
\dcref{apph:H.41}\group{Integral curves}

If $\gamma$ is locally Lipschitz at $s$, then its right tangent, when it exists, equals
the limit of the initial velocities of minimizing geodesics from $\gamma(s)$ to
$\gamma(s+\sigma)$, divided by $\sigma$.
\end{lemma}

\begin{remark}[Uniqueness of the right tangent]
\label{rem:chow-iii-h42-right-tangent-uniqueness}
\source{chow-et-al-ricci-flow-part-iii:page-0464}
\dcref{apph:H.42}\group{Integral curves}
The defining identity against all smooth functions makes the right tangent unique;
the exponential inverse is well-defined for sufficiently small increments.
\end{remark}

\begin{definition}[Integral curve of the normalized generalized gradient]
\label{def:chow-iii-h43-integral-curve}
\source{chow-et-al-ricci-flow-part-iii:page-0465}
\dcref{apph:H.43}\group{Integral curves}
For $f\in\Vzero(C)$ and $f(x)=a$, an integral curve is a locally Lipschitz path
$\gamma_x:[a,b]\to C$ with $\gamma_x(a)=x$, $f(\gamma_x(s))<f_{\sup}$, and
\[
(\dot\gamma_x)_+(s)=\frac{\nabla f(\gamma_x(s))}{|\nabla f(\gamma_x(s))|^2}.
\]
It is maximal when $f(b)=f_{\sup}$.
\end{definition}

\begin{example}[Integral curves in the square]
\label{ex:chow-iii-h44-square-integral-curves}
\source{chow-et-al-ricci-flow-part-iii:page-0465}
\dcref{apph:H.44}\group{Integral curves}
For $C=[-1,1]^2$ and $f=d(\cdot,\partial C)$, the maximal normalized-gradient
integral curve from a point runs linearly toward the central ridge, and thereafter
along that ridge to the maximum set, with $f$ itself as parameter.
\end{example}

\begin{lemma}[The function $f$ along an integral curve]
\label{lem:chow-iii-h45-integral-curve-level}
\source{chow-et-al-ricci-flow-part-iii:page-0465,chow-et-al-ricci-flow-part-iii:page-0466}
\uses{def:chow-iii-h43-integral-curve, lem:chow-iii-h41-right-tangent-geodesics}
\dcref{apph:H.45}\group{Integral curves}

If $\gamma_x:[a,b]\to C$ is an integral curve, then
\[
f(\gamma_x(s))=s\qquad(a\le s\le b).
\]
\end{lemma}

\begin{remark}[Model nonsmooth example]
\label{rem:chow-iii-h47-model-example}
\source{chow-et-al-ricci-flow-part-iii:page-0467}
\dcref{apph:H.47}\group{Integral curves}
The square example is a model for the nonsmooth generalized-gradient behavior.
\end{remark}

\begin{lemma}[Concavity after arclength reparametrization]
\label{lem:chow-iii-h48-integral-curve-concavity}
\source{chow-et-al-ricci-flow-part-iii:page-0467}
\uses{lem:chow-iii-h45-integral-curve-level, lem:chow-iii-h37-directional-gradient-inequality}
\dcref{apph:H.48}\group{Integral curves}

If an integral curve is reparametrized by arclength, then $f$ composed with the
reparametrized curve is concave.
\end{lemma}

\begin{lemma}[Monotonicity between equal-level integral curves]
\label{lem:chow-iii-h49-integral-curves-distance}
\source{chow-et-al-ricci-flow-part-iii:page-0467,chow-et-al-ricci-flow-part-iii:page-0468}
\uses{lem:chow-iii-h15-minimal-geodesics-locally-convex, lem:chow-iii-h38-acute-angle}
\dcref{apph:H.49}\group{Integral curves}

If two integral curves start at the same level, then
\[
s\longmapsto d_C(\gamma_x(s),\gamma_y(s))
\]
is nonincreasing.
\end{lemma}

\begin{corollary}[Uniqueness of integral curves]
\label{cor:chow-iii-h50-integral-curve-uniqueness}
\source{chow-et-al-ricci-flow-part-iii:page-0469}
\uses{lem:chow-iii-h49-integral-curves-distance}
\dcref{apph:H.50}\group{Integral curves}

Through each $x$ with $f(x)<f_{\sup}$ there is at most one integral curve of
$\nabla f/|\nabla f|^2$.
\end{corollary}

\begin{remark}[Square example and distance monotonicity]
\label{rem:chow-iii-h51-square-distance-exercise}
\source{chow-et-al-ricci-flow-part-iii:page-0469}
\dcref{apph:H.51}\group{Integral curves}
The equal-level distance monotonicity can be checked directly in the square example.
\end{remark}

\begin{lemma}[Distance to a higher level]
\label{lem:chow-iii-h52-integral-curve-to-point-distance}
\source{chow-et-al-ricci-flow-part-iii:page-0469}
\uses{lem:chow-iii-h45-integral-curve-level, lem:chow-iii-h38-acute-angle}
\dcref{apph:H.52}\group{Integral curves}

If $f(y)=b>a=f(x)$ and $\gamma_x:[a,b]\to C$ is an integral curve, then
\[
s\longmapsto d_C(\gamma_x(s),y)
\]
is nonincreasing.
\end{lemma}

\begin{proposition}[Existence of maximal integral curves]
\label{prop:chow-iii-h53-maximal-integral-curve}
\source{chow-et-al-ricci-flow-part-iii:page-0469,chow-et-al-ricci-flow-part-iii:page-0470}
\uses{lem:chow-iii-h52-integral-curve-to-point-distance, lem:chow-iii-h58-integral-curve-existence}
\dcref{apph:H.53}\group{Integral curves}

If $C$ is compact and $f\in\Vzero(C)$, every $x$ with $f(x)<f_{\sup}$ has a maximal
integral curve $\gamma_x:[f(x),f_{\sup}]\to C$.
\end{proposition}

\begin{lemma}[Distance to a superlevel set]
\label{lem:chow-iii-h54-distance-superlevel}
\source{chow-et-al-ricci-flow-part-iii:page-0470,chow-et-al-ricci-flow-part-iii:page-0471,chow-et-al-ricci-flow-part-iii:page-0472}
\uses{lem:chow-iii-h36-unique-generalized-gradient}
\dcref{apph:H.54}\group{Integral curves}

For $f(y_0)<f_{\sup}$ and $\epsilon>0$, there are a neighborhood $\mathcal U$ of
$y_0$ and $s_\epsilon>0$ such that
\[
\frac{s}{d_C(y,y^s)}\ge |\nabla f(y_0)|-\epsilon
\]
for $y\in C\cap\mathcal U$ and $0<s<s_\epsilon$, where $y^s$ is a nearest point in
the superlevel set $\{f\ge f(y)+s\}$.
\end{lemma}

\begin{lemma}[Uniform Lipschitz bound for broken geodesics]
\label{lem:chow-iii-h55-broken-geodesic-bound}
\source{chow-et-al-ricci-flow-part-iii:page-0472,chow-et-al-ricci-flow-part-iii:page-0473}
\dcref{apph:H.55}\group{Integral curves}
Every broken geodesic $\Gamma(x,P)$ built from nearest superlevel points satisfies
\[
d_C(\Gamma(x,P)(s),\Gamma(x,P)(s'))\le C|s-s'|,
\]
where $C$ depends only on $\operatorname{diam}(C)$, $f_{\sup}$, and the terminal level
below $f_{\sup}$, not on the partition $P$.
\end{lemma}

\begin{corollary}[Limits of broken geodesics]
\label{cor:chow-iii-h56-broken-geodesic-limit}
\source{chow-et-al-ricci-flow-part-iii:page-0473}
\uses{lem:chow-iii-h55-broken-geodesic-bound}
\dcref{apph:H.56}\group{Integral curves}

If the mesh of partitions $P_k$ tends to zero, a subsequence of $\Gamma(x,P_k)$
converges uniformly on every compact level interval to a path $\Gamma_\infty$.
\end{corollary}

\begin{lemma}[Properties of the limiting path]
\label{lem:chow-iii-h57-limit-path-properties}
\source{chow-et-al-ricci-flow-part-iii:page-0473,chow-et-al-ricci-flow-part-iii:page-0474}
\uses{lem:chow-iii-h54-distance-superlevel, lem:chow-iii-h55-broken-geodesic-bound}
\dcref{apph:H.57}\group{Integral curves}

The limiting path satisfies $\Gamma_\infty(f(x))=x$, is Lipschitz,
$f(\Gamma_\infty(s))=s$, and
\[
\limsup_{\sigma\to0^+}\frac{d_C(\Gamma_\infty(s+\sigma),\Gamma_\infty(s))}{\sigma}
\le |\nabla f(\Gamma_\infty(s))|^{-1}.
\]
\end{lemma}

\begin{lemma}[Existence of an integral curve below the top level]
\label{lem:chow-iii-h58-integral-curve-existence}
\source{chow-et-al-ricci-flow-part-iii:page-0474,chow-et-al-ricci-flow-part-iii:page-0475,chow-et-al-ricci-flow-part-iii:page-0476}
\uses{lem:chow-iii-h41-right-tangent-geodesics, lem:chow-iii-h57-limit-path-properties, lem:chow-iii-h37-directional-gradient-inequality}
\dcref{apph:H.58}\group{Integral curves}

Every uniform limit $\Gamma_\infty$ of vanishing-mesh broken geodesics is the
integral curve of $\nabla f/|\nabla f|^2$ emanating from $x$.
\end{lemma}

\begin{theorem}[Sharafutdinov distance-nonincreasing retraction]
\label{thm:chow-iii-h59-sharafutdinov-retraction}
\source{chow-et-al-ricci-flow-part-iii:page-0476,chow-et-al-ricci-flow-part-iii:page-0477}
\uses{prop:chow-iii-h53-maximal-integral-curve, cor:chow-iii-h50-integral-curve-uniqueness, lem:chow-iii-h52-integral-curve-to-point-distance, lem:chow-iii-h49-integral-curves-distance}
\dcref{apph:H.59}\group{Integral curves}

For $f\in\Vzero(C)$ and $C_s=\{f\ge s\}$, define
\[
\Gamma_s(x)=
\begin{cases}x,&f(x)\ge s,\\ \gamma_x(s),&f(x)<s.\end{cases}
\]
Then $\Gamma_s:C\to C_s$ is a retraction, is distance nonincreasing on $C$,
and the map $(x,s)\mapsto\Gamma_s(x)$ is continuous.  Consequently every rectifiable
path satisfies $\operatorname{L}(\Gamma_s\circ\alpha)\le\operatorname{L}(\alpha)$.
\end{theorem}

\begin{remark}[Global versus local Lipschitz constants]
\label{rem:chow-iii-h1-lipschitz-question}
\dcref{apph:H.1}\group{Euclidean convexity}\source{chow-et-al-ricci-flow-part-iii:page-0435}
For a map between metric spaces one may ask when
\[
\Lip(f)=\sup_{x}\dil_x(f).
\]
The equality is a global-geometric question; it holds for the real line and the
intrinsic circle metric, but need not follow from local Lipschitz regularity alone.
\end{remark}

\begin{remark}[Right continuity problem]
\label{rem:chow-iii-h46-right-continuity-problem}
\dcref{apph:H.46}\group{Integral curves}\source{chow-et-al-ricci-flow-part-iii:page-0466}
The right continuity of $\nabla f$ and $\nabla f/|\nabla f|^2$ along their integral curves
is the natural regularity question behind the level identity above.
\end{remark}
